\documentclass[10pt,a4paper]{amsart}
\usepackage{amsmath,amssymb,amsthm,graphicx}
% Retained minimal preamble for the Stacks Project chapter more-algebra.tex (More on Algebra).
% The chapter file pinned-original-source.tex carries only \input{preamble} on its first line, so no
% preamble text lives inside the chapter itself.  The Stacks Project uses one single shared preamble for
% every chapter; the snapshot stages it identically as sources/stacks-chow/preamble.tex and
% sources/stacks-etale-cohomology/preamble.tex.  The block below is copied mechanically, line for line,
% from sources/stacks-chow/preamble.tex lines 14-18 (the Xy-pic setup) and 154-217 (the native theorem
% declarations together with the mathematical macros).  Nothing was added, removed or redefined; the
% wrapper class and the amsmath/amssymb/amsthm/graphicx base packages are supplied by the assembler.
% Chapter: more-algebra.tex -- More on Algebra.
% For commutative diagrams we use Xy-pic
\usepackage[all]{xy}

% We use 2cell for 2-commutative diagrams.
\xyoption{2cell}

% Theorem environments.
%
\theoremstyle{plain}
\newtheorem{theorem}[subsection]{Theorem}
\newtheorem{proposition}[subsection]{Proposition}
\newtheorem{lemma}[subsection]{Lemma}

\theoremstyle{definition}
\newtheorem{definition}[subsection]{Definition}
\newtheorem{example}[subsection]{Example}
\newtheorem{exercise}[subsection]{Exercise}
\newtheorem{situation}[subsection]{Situation}

\theoremstyle{remark}
\newtheorem{remark}[subsection]{Remark}
\newtheorem{remarks}[subsection]{Remarks}

\numberwithin{equation}{subsection}

% Macros
%
\def\lim{\mathop{\mathrm{lim}}\nolimits}
\def\colim{\mathop{\mathrm{colim}}\nolimits}
\def\Spec{\mathop{\mathrm{Spec}}}
\def\Hom{\mathop{\mathrm{Hom}}\nolimits}
\def\Ext{\mathop{\mathrm{Ext}}\nolimits}
\def\SheafHom{\mathop{\mathcal{H}\!\mathit{om}}\nolimits}
\def\SheafExt{\mathop{\mathcal{E}\!\mathit{xt}}\nolimits}
\def\Sch{\mathit{Sch}}
\def\Mor{\mathop{\mathrm{Mor}}\nolimits}
\def\Ob{\mathop{\mathrm{Ob}}\nolimits}
\def\Sh{\mathop{\mathit{Sh}}\nolimits}
\def\NL{\mathop{N\!L}\nolimits}
\def\CH{\mathop{\mathrm{CH}}\nolimits}
\def\proetale{{pro\text{-}\acute{e}tale}}
\def\etale{{\acute{e}tale}}
\def\QCoh{\mathit{QCoh}}
\def\Ker{\mathop{\mathrm{Ker}}}
\def\Im{\mathop{\mathrm{Im}}}
\def\Coker{\mathop{\mathrm{Coker}}}
\def\Coim{\mathop{\mathrm{Coim}}}

% Boxtimes
%
\DeclareMathSymbol{\boxtimes}{\mathbin}{AMSa}{"02}

%
% Macros for moduli stacks/spaces
%
\def\QCohstack{\mathcal{QC}\!\mathit{oh}}
\def\Cohstack{\mathcal{C}\!\mathit{oh}}
\def\Spacesstack{\mathcal{S}\!\mathit{paces}}
\def\Quotfunctor{\mathrm{Quot}}
\def\Hilbfunctor{\mathrm{Hilb}}
\def\Curvesstack{\mathcal{C}\!\mathit{urves}}
\def\Polarizedstack{\mathcal{P}\!\mathit{olarized}}
\def\Complexesstack{\mathcal{C}\!\mathit{omplexes}}
% \Pic is the operator that assigns to X its picard group, usage \Pic(X)
% \Picardstack_{X/B} denotes the Picard stack of X over B
% \Picardfunctor_{X/B} denotes the Picard functor of X over B
\def\Pic{\mathop{\mathrm{Pic}}\nolimits}
\def\Picardstack{\mathcal{P}\!\mathit{ic}}
\def\Picardfunctor{\mathrm{Pic}}
\def\Deformationcategory{\mathcal{D}\!\mathit{ef}}
\title{Lifting an inverse system of objects of D(Ab) to a system of complexes whose transition diagrams commute strictly}
\author{The Stacks Project Authors}
\date{stacks-project:more-algebra.tex:local-snapshot-20261012; GNU FDL 1.2 or later}
\begin{document}
\maketitle
\setcounter{subsection}{1}
\noindent\textit{Editorial note.} All delivered mathematical text is the author's own, copied from the pinned Stacks Project chapter file more-algebra.tex; no statement or proof step was rewritten, reordered or removed. The counted statement is the labeled lemma environment at source lines 23334-23348 (canonical tag 0CQ9) and the counted proof is the proof environment at source lines 23350-23373, which immediately follows it and is its own proof.\par\medskip
\noindent\textit{Editorial note.} Both the statement and the proof display a diagram in the derived category with Xy-pic, using the author's own markup and the declaration of Ab and of N from the retained preamble; nothing about the diagrams was changed.\par\medskip
\noindent\textit{Editorial note.} The lead-in text at source lines 23330-23332, contrasting an inverse system of complexes with an inverse system of objects of D(Ab) and announcing that the next lemma and remark explain the difference, is outside the counted spans because it belongs to the surrounding discussion.\par\medskip
\noindent\textit{Editorial note.} The chapter file pinned-original-source.tex carries only an input of the shared preamble on its first line, so no preamble text lives inside the chapter. The retained file sources/stacks-more-algebra/retained-preamble.tex is a mechanical, line-for-line copy of the project's single shared preamble as staged in the snapshot (Xy-pic setup plus the native theorem declarations and mathematical macros). Nothing was added, removed or redefined, and the excerpt serves only as preamble, never as a span coordinate.\par\medskip
\noindent\textit{Editorial note.} The source URL is the source-level Stacks Project URL, as required when more than one card is produced from one source; the per-claim tag is carried by claim\_key and by this note. The counted statement and proof refer to no labelled item of the chapter and cite no bibliography entry, so this card contributes no external reference stub and needs no bibliography.\par\medskip
\noindent\textit{Allowed standard prerequisites (editorial).}\begin{itemize}
\item The objects of D(Ab) are complexes of abelian groups up to homotopy of maps, and a morphism in D(Ab) between two complexes is represented by a morphism of complexes after resolving the source by a complex of projectives. used at proof lines 23363-23367, where the author appeals to the definition of morphisms in D(Ab) to find a complex and a morphism of complexes representing a prescribed composition. Supplies the transition map of complexes that represents the composite in D(Ab), which is what makes the resulting system a system of complexes rather than a system of derived objects.
\item The category D(Ab(N)) of inverse systems of complexes of abelian groups has as morphisms the levelwise morphisms of complexes, and two such systems are isomorphic when every level is an isomorphism in D(Ab) and the transition maps agree in D(Ab). used at proof lines 23336-23347, where the object M = (M\_n) of D(Ab(N)) and the isomorphisms M\_n -> K\_n are required. Fixes the exact form of the lifting that has to be constructed level by level.
\item An inverse system indexed by the natural numbers may be constructed by induction on the length of its finite initial segments. used at proof lines 23351-23355 and closed at proof line 23372. Turns the levelwise construction into an induction, with the case of length one as the initial step.
\end{itemize}\medskip
\begin{lemma}
\label{lemma-lift-to-system-complexes-Ab}
Let $(K_n)$ be an inverse system of objects of $D(\textit{Ab})$.
Then there exists an object $M = (M_n^\bullet)$
of $D(\textit{Ab}(\mathbf{N}))$ and isomorphisms
$M_n^\bullet \to K_n$ in $D(\textit{Ab})$ such that the diagrams
$$
\xymatrix{
M_{n + 1}^\bullet \ar[d] \ar[r] &
M_n^\bullet \ar[d] \\
K_{n + 1} \ar[r] & K_n
}
$$
commute in $D(\textit{Ab})$.
\end{lemma}

\begin{proof}
Namely, let $M_1^\bullet$ be a complex of abelian groups representing $K_1$.
Suppose we have constructed
$M_e^\bullet \to M_{e - 1}^\bullet \to \ldots \to M_1^\bullet$
and maps $\psi_i : M_i^\bullet \to K_i$ such that
the diagrams in the statement of the lemma commute for all $n < e$.
Then we consider the diagram
$$
\xymatrix{
& M_n^\bullet \ar[d]^{\psi_n} \\
K_{n + 1} \ar[r] & K_n
}
$$
in $D(\textit{Ab})$. By the definition of morphisms in $D(\textit{Ab})$
we can find a complex $M_{n + 1}^\bullet$ of abelian groups, an isomorphism
$M_{n + 1}^\bullet \to K_{n + 1}$ in $D(\textit{Ab})$, and a
morphism of complexes $M_{n + 1}^\bullet \to M_n^\bullet$
representing the composition
$$
K_{n + 1} \to K_n \xrightarrow{\psi_n^{-1}} M_n^\bullet
$$
in $D(\textit{Ab})$.
Thus the lemma holds by induction.
\end{proof}

\end{document}
